% answers_short.tex —— 2013 年·浙江卷·理科 · 简短答案速查
% 由 tools/make_short_answers.py 自 answers.tex 的 \gksolhead/\gkanswer 提取，
% 并与 2013/2013_数学真题_浙江理_答案要点.md 逐题比对后生成，勿手改（改源 MD 后重跑）。

\begin{gknotice}
\textbf{使用说明}：本节只给最终答案便于快速核对，完整推导见解析版。答案为编写者独立推导，\textbf{非官方标准答案}。
\end{gknotice}

\gkpart{选择题}
\begin{center}\begin{tabular}{ccccccccccc}
\toprule
题号 & 1 & 2 & 3 & 4 & 5 & 6 & 7 & 8 & 9 & 10 \\
\midrule
答案 & B & C & D & B & A & C & D & C & D & A \\
\bottomrule
\end{tabular}\end{center}

\gkpart{填空题}
\begin{itemize}[leftmargin=2.6em]
  \item[\textbf{11.}] $-10$。通项为 $\mathrm C_5^r(-1)^rx^{\frac{15-5r}6}$，令 $15-5r=0$ 得 $r=3$，故 $A=-\mathrm C_5^3=-10$。
  \item[\textbf{12.}] $24$。该几何体是底面为两直角边 $3$、$4$ 的直角三角形、高为 $5$ 的直三棱柱截去一个三棱锥所得，$V=\dfrac12\cdot3\cdot4\cdot5-\dfrac13\cdot\dfrac12\cdot3\cdot4\cdot3=30-6=24$。
  \item[\textbf{13.}] $2$。可行域顶点为 $(0,2)$，$(2,0)$，$(4,4)$，对应 $z$ 值为 $2$，$2k$，$4k+4$；由最大值 $12$ 得 $k=2$。
  \item[\textbf{14.}] $480$。六个字母共 $6!=720$ 种排法，其中 $C$ 在 $A$、$B$ 之间占 $\dfrac26$，故 $A$、$B$ 同侧的排法为 $720\times\dfrac46=480$。
  \item[\textbf{15.}] $\pm1$。设 $A$、$B$ 的横坐标之和为 $2x_0$，由中点弦方程及过 $P(-1,0)$ 得 $x_0=1$，再由 $(x_0-1)^2+y_0^2=4$ 得 $y_0=\pm2$，斜率 $k=\dfrac2{y_0}=\pm1$。
  \item[\textbf{16.}] $\dfrac{\sqrt6}3$。设 $C(0,0)$，$A(b,0)$，$B(0,a)$，由 $\sin\angle BAM=\dfrac13$ 得 $b^2=\dfrac{a^2}2$，故 $\sin\angle BAC=\dfrac{a}{\sqrt{a^2+b^2}}=\sqrt{\dfrac23}=\dfrac{\sqrt6}3$。
  \item[\textbf{17.}] $2$。$|\boldsymbol b|^2=x^2+y^2+\sqrt3xy=x^2\left[\left(\dfrac yx\right)^2+\sqrt3\cdot\dfrac yx+1\right]\geqslant\dfrac{x^2}4$，故 $\dfrac{|x|}{|\boldsymbol b|}\leqslant2$。
\end{itemize}

\gkpart{解答题}
\begin{itemize}[leftmargin=2.6em]
  \item[\textbf{18.}] (1) $d=4$，$a_n=4n+6$；或 $d=-1$，$a_n=11-n$。由 $(2a_2+2)^2=5a_1a_3$ 得 $d^2-3d-4=0$；(2) 当 $n\leqslant11$ 时，$\dfrac{n(21-n)}2$；当 $n\geqslant12$ 时，$\dfrac{n^2-21n+220}2$。此时 $a_n=11-n$，前 $11$ 项绝对值和为 $55$，其后逐项增大再累加
  \item[\textbf{19.}] (1) $P(\xi=2)=\dfrac14$，$P(\xi=3)=\dfrac13$，$P(\xi=4)=\dfrac5{18}$，$P(\xi=5)=\dfrac19$，$P(\xi=6)=\dfrac1{36}$。单次得分为 $1,2,3$ 的概率为 $\dfrac12,\dfrac13,\dfrac16$，两次独立相加；(2) $a:b:c=3:2:1$。由 $E\eta=\dfrac53$、$E(\eta^2)=D\eta+(E\eta)^2=\dfrac{10}3$ 解得 $a=3c$，$b=2c$
  \item[\textbf{20.}] (1) 证明见解析。取 $D$ 为原点、$DC$、$DA$ 为 $x$、$z$ 轴建系，则 $\overrightarrow{PQ}$ 的竖坐标为 $0$，故 $PQ\parallel$ 平面 $BCD$；(2) $60^\circ$。由二面角为 $60^\circ$ 解得 $CD^2=2$，故 $\cos\angle BDC=\dfrac{CD}{BD}=\dfrac12$
  \item[\textbf{21.}] (1) $\dfrac{x^2}4+y^2=1$。由顶点 $P(0,-1)$ 得 $b=1$，由长轴是圆 $x^2+y^2=4$ 的直径得 $a=2$；(2) $y=\dfrac{\sqrt{10}}2x-1$ 或 $y=-\dfrac{\sqrt{10}}2x-1$。设 $l_1:y=kx-1$，则 $S_{\triangle ABD}=\dfrac{8\sqrt{4k^2+3}}{k^2+4}$，令 $u=4k^2+3$ 知 $u=13$ 时取最大值，即 $k^2=\dfrac52$
  \item[\textbf{22.}] (1) $y=(3a-3)x+4-3a$。$f(1)=1$，$f'(1)=3a-3$；(2) 当 $a\leqslant0$ 时最大值为 $3-3a$；当 $0<a\leqslant\dfrac34$ 时最大值为 $1+2(1-a)^{\frac32}$；当 $a>\dfrac34$ 时最大值为 $3a-1$。令 $t=x-1$，问题化为 $\left|t^3+3(a-1)t+1\right|$ 在 $[-1,1]$ 上的最大值
\end{itemize}
